\subsubsection{Banach 空间的例子}\label{subsec:banach-space-examples}

\Needspace{4\baselineskip}
\begin{Example}
\label{ex:bounded-continuous-functions-banach}
设 $\Omega$ 为 $\mathbb{R}^m$ 中的开集, 用 $C_b(\overline\Omega)$ 表示 $\overline\Omega$ 上所有有界连续函数的集合. 定义
\begin{equation}\label{eq:bounded-continuous-functions-supremum-norm}
\|u\|=\sup_{\bm{x}\in\overline\Omega}|u(\bm{x})|.
\end{equation}
那么 $C_b(\overline\Omega)$ 是 Banach 空间. 事实上, 范数的性质容易检验, 而完备性就是连续函数的一致极限连续这一熟知事实的重述 (注意, $C_b(\overline\Omega)$ 中的收敛就是一致收敛).
\end{Example}

\Needspace{4\baselineskip}
\begin{Example}
\label{ex:bounded-analytic-functions-banach}
设 $\Omega$ 为 $\mathbb{C}$ 中的开集, 用 $A(\Omega)$ 表示 $\Omega$ 上所有有界解析函数的集合. 定义
\begin{equation}\label{eq:bounded-analytic-functions-supremum-norm}
\|u\|=\sup_{z\in\Omega}|u(z)|.
\end{equation}
那么 $A(\Omega)$ 是 Banach 空间. 事实上, 解析函数的一致极限仍解析这一事实, 是 Morera 定理的一个容易得到的推论.
\end{Example}

\begin{Example}
\label{ex:continuous-functions-lp-norm}
设 $\Omega$ 为 $\mathbb{R}^m$ 中的局部 Jordan 可测开集 (参见 \cite{Bu}) (这意味着它与任意球的交均 Jordan 可测), 并设 $1\leq p<\infty$. 用 $\widetilde L^p(\Omega)$ 表示所有使
\begin{equation}\label{eq:continuous-functions-lp-integral-norm}
\|u\|_p:=\left(\int_\Omega|u(\bm{x})|^p\,\mathrm{d}\bm{x}\right)^{1/p}
\end{equation}
有限的连续函数 $u:\overline\Omega\to\mathbb{R}$ 的集合. 要证明 $\widetilde L^p(\Omega)$ 为赋范向量空间. 范数的前两条性质显然成立, 而当 $p=1$ 时三角不等式也显然成立. 当 $p>1$ 时, 三角不等式需要一些准备.
\end{Example}

\begin{Lemma}
\label{lem:young-conjugate-exponents-scalar-inequality}
设 $a,b\geq0$, $p,q>1$ 且 $p^{-1}+q^{-1}=1$. 那么
\begin{equation}\label{eq:young-conjugate-exponents-scalar-inequality}
ab\leq\dfrac{a^p}p+\dfrac{b^q}q.
\end{equation}

\end{Lemma}

\PDEIndexMark{idx:young-inequality-minimum-argument}对函数 \[b\mapsto a^p/p+b^q/q-ab.\] 求最小值即可得到证明. 见 Problem~\ref{ex:young-conjugate-exponents-inequality}. 现在由 Lemma~\ref{lem:young-conjugate-exponents-scalar-inequality} 建立 Hölder 不等式.

\begin{Theorem}
\label{thm:continuous-functions-holder-inequality}
设 $f\in\widetilde L^p(\Omega)$, $g\in\widetilde L^q(\Omega)$, 其中 $p,q\in(1,\infty)$ 且 $p^{-1}+q^{-1}=1$. 那么 $fg\in\widetilde L^1(\Omega)$, 并且
\begin{equation}\label{eq:continuous-functions-holder-inequality}
\|fg\|_1\leq\|f\|_p\|g\|_q.
\end{equation}

\end{Theorem}

\begin{Proof}
只需假设 $\|f\|_p\ne0$, $\|g\|_q\ne0$ 以排除平凡情形. 由前一引理, 得到
\begin{equation}\label{eq:holder-normalized-young-pointwise-bound}
\dfrac{|f(\bm{x})|}{\|f\|_p}\dfrac{|g(\bm{x})|}{\|g\|_q}\leq\dfrac1p\dfrac{|f(\bm{x})|^p}{\|f\|_p^p}+\dfrac1q\dfrac{|g(\bm{x})|^q}{\|g\|_q^q}.
\end{equation}
在 $\Omega$ 上积分即得该定理.
\end{Proof}

\begin{Theorem}
\label{thm:continuous-functions-lp-triangle-inequality}
设 $f,g\in\widetilde L^p(\Omega)$. 那么 $f+g\in\widetilde L^p(\Omega)$ 且 $\|f+g\|_p\leq\|f\|_p+\|g\|_p$.
\end{Theorem}

\begin{Proof}
注意
\begin{equation}\label{eq:lp-sum-integrability-bound}
|f+g|^p\leq(|f|+|g|)^p\leq2^p(|f|^p+|g|^p);
\end{equation}
因而 $f+g\in\widetilde L^p(\Omega)$. 此外, 有
\begin{equation}\label{eq:lp-triangle-pointwise-product-bound}
|f+g|^p\leq|f+g|^{p-1}|f|+|f+g|^{p-1}|g|,
\end{equation}
由 $p^{-1}+q^{-1}=1$, 有 $(p-1)q=p$; 因而 $|f+g|^{p-1}\in\widetilde L^q(\Omega)$. 因此可应用 Hölder 不等式, 得到
\begin{equation}\label{eq:lp-triangle-holder-integral-bound}
\int_\Omega|f+g|^p\,\mathrm{d}\bm{x}\leq\left(\int_\Omega|f+g|^p\,\mathrm{d}\bm{x}\right)^{1/q}(\|f\|_p+\|g\|_p).
\end{equation}
由此立即得出三角不等式.
\end{Proof}
